\documentclass[11pt]{article}
\usepackage{amsmath,amssymb,amsthm}
\usepackage{algorithm}
\usepackage[noend]{algpseudocode}
\usepackage{geometry}
\geometry{margin=1in}
\newtheorem{theorem}{Theorem}
\newtheorem{lemma}{Lemma}
\newtheorem{assumption}{Assumption}
\newcommand{\E}{\mathbb{E}}
\newcommand{\R}{\mathbb{R}}

\begin{document}

\title{AdaGrad‑Style Gradient Averaging for Non‑Convex Smooth Optimization}
\author{Compiled \\ \today}
\maketitle

%%%%%%%%%%%%%%%%%%%%%%%%%%%%%%%%%%%%%%%%%%%%%%%%%%%%%%%%%%%%%%%%%%%%%%%%%%%%
\section*{Algorithm Name}
\textbf{AdaGrad with Running Gradient Sum (RG‑AdaGrad)}  

The method keeps a cumulative sum of past gradients and adapts the stepsize
according to the accumulated squared magnitude of those gradients.

%%%%%%%%%%%%%%%%%%%%%%%%%%%%%%%%%%%%%%%%%%%%%%%%%%%%%%%%%%%%%%%%%%%%%%%%%%%%
\section*{Mathematical Formulation}
Let $f:\R^{d}\rightarrow\R$ be differentiable.  
Denote the stochastic (or full) gradient at iteration $k$ by 
\[
g_{k}=\nabla f(w_{k})\in\R^{d}.
\]

Define the cumulative (vector) sum of gradients and of squared gradients:
\[
A_{k}\;:=\;\sum_{i=1}^{k} g_{i},\qquad
G_{k}\;:=\;\sum_{i=1}^{k} g_{i}\odot g_{i}\;\in\R^{d},
\]
where $\odot$ denotes element‑wise product.  
For a scalar hyper‑parameter $\alpha>0$ and a small constant $\varepsilon>0$,
the adaptive stepsize vector is
\[
\eta_{k}\;:=\;\frac{\alpha}{\sqrt{G_{k}+\varepsilon\,\mathbf 1}}\;\in\R^{d},
\]
(the square root and division are element‑wise).  
The update rule is
\[
\boxed{%
w_{k+1}=w_{k}-\eta_{k}\odot g_{k}
}\tag{1}
\]
with the initialisation $w_{0}\in\R^{d}$, $A_{0}=0$, $G_{0}=0$.

For the one‑dimensional case (used in the dry‑run) the formulas reduce to
\[
A_{k}= \sum_{i=1}^{k} g_{i},\qquad
G_{k}= \sum_{i=1}^{k} g_{i}^{2},\qquad
\eta_{k}= \frac{\alpha}{\sqrt{G_{k}+\varepsilon}} ,
\qquad
w_{k+1}= w_{k}-\eta_{k} g_{k}.
\]

%%%%%%%%%%%%%%%%%%%%%%%%%%%%%%%%%%%%%%%%%%%%%%%%%%%%%%%%%%%%%%%%%%%%%%%%%%%%
\section*{Pseudocode}
\begin{algorithm}[H]
\caption{RG‑AdaGrad}
\begin{algorithmic}[1]
\Require Initial point $w_{0}\in\R^{d}$, stepsize base $\alpha>0$, smoothing $\varepsilon>0$.
\State $A_{0}\gets 0$, $G_{0}\gets 0$.
\For{$k=0,1,\dots,T-1$}
    \State Compute (stochastic) gradient $g_{k}\gets \nabla f(w_{k})$.
    \State $A_{k+1}\gets A_{k}+g_{k}$.
    \State $G_{k+1}\gets G_{k}+g_{k}\odot g_{k}$.
    \State $\eta_{k}\gets \displaystyle\frac{\alpha}{\sqrt{G_{k+1}+\varepsilon\,\mathbf 1}}$.
    \State $w_{k+1}\gets w_{k}-\eta_{k}\odot g_{k}$.
\EndFor
\State \textbf{return} $w_{T}$ (or any averaged iterate).
\end{algorithmic}
\end{algorithm}

%%%%%%%%%%%%%%%%%%%%%%%%%%%%%%%%%%%%%%%%%%%%%%%%%%%%%%%%%%%%%%%%%%%%%%%%%%%%
\section*{Dry Run Example}
Consider the scalar quadratic $f(w)=(w-3)^{2}$, gradient $g(w)=2(w-3)$.
Set $w_{0}=0$, $\alpha=0.1$, $\varepsilon=10^{-8}$ and run three iterations.

\begin{center}
\begin{tabular}{c|c|c|c|c}
$k$ & $w_{k}$ & $g_{k}=2(w_{k}-3)$ & $G_{k}=\sum_{i=1}^{k}g_{i}^{2}$ & $w_{k+1}=w_{k}-\dfrac{\alpha}{\sqrt{G_{k}+\varepsilon}}g_{k}$\\\hline
0 & $0$ & $-6$ & $36$ &
$w_{1}=0-\dfrac{0.1}{\sqrt{36}}(-6)=0+0.1=0.1000$\\[4pt]
1 & $0.1000$ & $-5.8$ & $36+33.64=69.64$ &
$w_{2}=0.1000-\dfrac{0.1}{\sqrt{69.64}}(-5.8)
          \approx0.1000+0.0694=0.1694$\\[4pt]
2 & $0.1694$ & $-5.6612$ &
$69.64+32.047=101.687$ &
$w_{3}=0.1694-\dfrac{0.1}{\sqrt{101.687}}(-5.6612)
          \approx0.1694+0.0560=0.2254$
\end{tabular}
\end{center}
Objective values:
\[
f(w_{0})=9,\; f(w_{1})\approx8.41,\; f(w_{2})\approx8.01,\; f(w_{3})\approx7.55.
\]

The example illustrates that the adaptive stepsize shrinks as the accumulated
squared gradient $G_{k}$ grows, leading to a more cautious update.

%%%%%%%%%%%%%%%%%%%%%%%%%%%%%%%%%%%%%%%%%%%%%%%%%%%%%%%%%%%%%%%%%%%%%%%%%%%%
\section*{Assumptions}
\begin{assumption}[Smoothness]\label{ass:smooth}
$f$ is $L$‑smooth, i.e.
\[
\forall x,y\in\R^{d}:\qquad 
f(y)\le f(x)+\langle\nabla f(x),y-x\rangle+\frac{L}{2}\|y-x\|^{2}.
\]
\end{assumption}
\begin{assumption}[Bounded Stochastic Gradient]\label{ass:bound}
There exists $G_{\max}>0$ such that for all $k$,
\[
\|g_{k}\|_{\infty}\le G_{\max}\quad\text{(element‑wise bound).}
\]
Consequently $0\le G_{k}\le k G_{\max}^{2}\mathbf 1$.
\end{assumption}
\begin{assumption}[Unbiased Gradient]\label{ass:unbiased}
The stochastic gradient satisfies
\[
\E[g_{k}\mid w_{k}]=\nabla f(w_{k}) .
\]
\end{assumption}
All expectations are taken w.r.t. the randomness of the stochastic oracle.

%%%%%%%%%%%%%%%%%%%%%%%%%%%%%%%%%%%%%%%%%%%%%%%%%%%%%%%%%%%%%%%%%%%%%%%%%%%%
\section*{Main Convergence Theorem}
\begin{theorem}\label{thm:main}
Assume \ref{ass:smooth}–\ref{ass:unbiased} hold and set $\varepsilon>0$.
Let $\{w_{k}\}_{k=0}^{T}$ be generated by RG‑AdaGrad with stepsize base
$\alpha>0$. Then for any $T\ge1$,
\[
\frac{1}{T}\sum_{k=0}^{T-1}\E\!\bigl[\|\nabla f(w_{k})\|^{2}\bigr]
\;\le\;
\frac{2\bigl(f(w_{0})-f^{\star}\bigr)}{\alpha\sqrt{T}}
\;+\;
\frac{L\alpha G_{\max}^{2}}{\sqrt{T}}
\;+\;
\frac{L\varepsilon d}{\alpha\sqrt{T}},
\tag{2}
\]
where $f^{\star}=\inf_{w}f(w)$ and $d$ is the dimension.
Consequently,
\[
\min_{0\le k<T}\E\!\bigl[\|\nabla f(w_{k})\|^{2}\bigr]
=O\!\bigl(T^{-1/2}\bigr).
\]
\end{theorem}

%%%%%%%%%%%%%%%%%%%%%%%%%%%%%%%%%%%%%%%%%%%%%%%%%%%%%%%%%%%%%%%%%%%%%%%%%%%%
\section*{Theoretical Properties}
\begin{itemize}
\item \textbf{Stability.} Because the denominator $\sqrt{G_{k}+\varepsilon}$ is non‑decreasing,
the effective stepsize $\eta_{k}$ never exceeds $\alpha/\sqrt{\varepsilon}$,
preventing exploding updates.
\item \textbf{Hyper‑parameter Sensitivity.}
  \begin{itemize}
    \item Larger $\alpha$ speeds up early progress (larger $\eta_{0}$) but may increase the second term in \eqref{thm:main}.
    \item $\varepsilon$ only influences an additive $O(\varepsilon d/\sqrt{T})$ term; setting $\varepsilon$ to $10^{-8}$ is standard.
  \end{itemize}
\item \textbf{Comparison to Baselines.}
  \begin{itemize}
    \item Fixed‑step SGD under the same assumptions yields an $O(T^{-1/2})$ bound on the average gradient norm but requires a manually tuned diminishing stepsize.
    \item RG‑AdaGrad automatically adapts the stepsize, offering a comparable rate without explicit schedule.
  \end{itemize}
\end{itemize}

%%%%%%%%%%%%%%%%%%%%%%%%%%%%%%%%%%%%%%%%%%%%%%%%%%%%%%%%%%%%%%%%%%%%%%%%%%%%
\section*{Discussion}
The bound \eqref{thm:main} mirrors the classic AdaGrad guarantee for convex
problems, yet it holds for general non‑convex $L$‑smooth objectives.
The $1/\sqrt{T}$ decay is optimal for first‑order stochastic methods under
the stated assumptions.  The proof below follows the smoothness inequality,
expands the adaptive stepsize, and carefully controls the stochastic noise
through the bounded‑gradient assumption.

%%%%%%%%%%%%%%%%%%%%%%%%%%%%%%%%%%%%%%%%%%%%%%%%%%%%%%%%%%%%%%%%%%%%%%%%%%%%
\section*{Preliminaries (Definitions and Lemmas)}
\begin{lemma}[Smoothness Descent Lemma]\label{lem:smooth}
If $f$ is $L$‑smooth then for any $x$ and any step $s$,
\[
f(x+s)\le f(x)+\langle\nabla f(x),s\rangle+\frac{L}{2}\|s\|^{2}.
\]
\end{lemma}
\begin{proof}
Standard; follows from integrating the Lipschitz condition on the gradient.
\end{proof}

\begin{lemma}[Bound on Cumulative Adaptive Stepsizes]\label{lem:stepbound}
For any $k\ge1$,
\[
\frac{\alpha^{2}}{2}\sum_{i=0}^{k-1}\frac{\|g_{i}\|^{2}}{G_{i+1}+\varepsilon}
\;\le\;
\frac{\alpha^{2}}{2\varepsilon}\sum_{i=0}^{k-1}\|g_{i}\|^{2}
\;\le\;
\frac{\alpha^{2}G_{\max}^{2}k}{2\varepsilon}.
\]
\end{lemma}
\begin{proof}
Since $G_{i+1}\ge 0$, $G_{i+1}+\varepsilon\ge\varepsilon$,
hence each fraction is bounded by $1/\varepsilon$.
\end{proof}

%%%%%%%%%%%%%%%%%%%%%%%%%%%%%%%%%%%%%%%%%%%%%%%%%%%%%%%%%%%%%%%%%%%%%%%%%%%%
\section*{Main Proof (Step‑by‑Step Derivation)}
We prove Theorem~\ref{thm:main}.  Throughout we condition on the filtration
$\mathcal{F}_{k}=\sigma(g_{0},\dots,g_{k-1})$.

\medskip
\noindent\textbf{Step 1. Apply smoothness.}  
From Lemma~\ref{lem:smooth} with $s=-\eta_{k}g_{k}$,
\[
f(w_{k+1})\le f(w_{k})-\eta_{k}\langle g_{k},g_{k}\rangle
+\frac{L}{2}\eta_{k}^{2}\|g_{k}\|^{2}.
\tag{3}
\]

\noindent\textbf{Step 2. Insert the definition of $\eta_{k}$.}  
Recall $\eta_{k}= \alpha/( \sqrt{G_{k+1}+\varepsilon})$, thus
\[
\eta_{k}\|g_{k}\|^{2}= \frac{\alpha\|g_{k}\|^{2}}{\sqrt{G_{k+1}+\varepsilon}},\qquad
\eta_{k}^{2}\|g_{k}\|^{2}= \frac{\alpha^{2}\|g_{k}\|^{2}}{G_{k+1}+\varepsilon}.
\tag{4}
\]

\noindent\textbf{Step 3. Rewrite (3) using (4).}
\[
f(w_{k+1})\le f(w_{k})
-\frac{\alpha\|g_{k}\|^{2}}{\sqrt{G_{k+1}+\varepsilon}}
+\frac{L\alpha^{2}}{2}\frac{\|g_{k}\|^{2}}{G_{k+1}+\varepsilon}.
\tag{5}
\]

\noindent\textbf{Step 4. Take conditional expectation.}  
Using Assumption~\ref{ass:unbiased},
\[
\E\!\bigl[f(w_{k+1})\mid\mathcal{F}_{k}\bigr]
\le f(w_{k})
-\frac{\alpha\|\nabla f(w_{k})\|^{2}}{\sqrt{G_{k+1}+\varepsilon}}
+\frac{L\alpha^{2}}{2}\frac{\E[\|g_{k}\|^{2}\mid\mathcal{F}_{k}]}{G_{k+1}+\varepsilon}.
\tag{6}
\]

\noindent\textbf{Step 5. Upper‑bound the variance term.}  
By Assumption~\ref{ass:bound}, $\|g_{k}\|_{\infty}\le G_{\max}$, hence
$\|g_{k}\|^{2}\le d\,G_{\max}^{2}$, and also $\|\nabla f(w_{k})\|^{2}\le
\|g_{k}\|^{2}$.  Therefore
\[
\E[\|g_{k}\|^{2}\mid\mathcal{F}_{k}]
\le d\,G_{\max}^{2}.
\tag{7}
\]

\noindent\textbf{Step 6. Substitute (7) into (6).}
\[
\E\!\bigl[f(w_{k+1})\mid\mathcal{F}_{k}\bigr]
\le f(w_{k})
-\frac{\alpha\|\nabla f(w_{k})\|^{2}}{\sqrt{G_{k+1}+\varepsilon}}
+\frac{L\alpha^{2}d\,G_{\max}^{2}}{2\bigl(G_{k+1}+\varepsilon\bigr)}.
\tag{8}
\]

\noindent\textbf{Step 7. Rearrange to isolate the gradient term.}
\[
\frac{\alpha\|\nabla f(w_{k})\|^{2}}{\sqrt{G_{k+1}+\varepsilon}}
\le f(w_{k})-\E\!\bigl[f(w_{k+1})\mid\mathcal{F}_{k}\bigr]
+\frac{L\alpha^{2}d\,G_{\max}^{2}}{2\bigl(G_{k+1}+\varepsilon\bigr)}.
\tag{9}
\]

\noindent\textbf{Step 8. Sum (9) over $k=0,\dots,T-1$ and take full expectation.}
Using telescoping of the first two terms,
\[
\alpha\sum_{k=0}^{T-1}
\E\!\left[\frac{\|\nabla f(w_{k})\|^{2}}{\sqrt{G_{k+1}+\varepsilon}}\right]
\le f(w_{0})- \E[f(w_{T})]
+\frac{L\alpha^{2}d\,G_{\max}^{2}}{2}
\sum_{k=0}^{T-1}\E\!\left[\frac{1}{G_{k+1}+\varepsilon}\right].
\tag{10}
\]

\noindent\textbf{Step 9. Lower bound $\sqrt{G_{k+1}+\varepsilon}$ by $\sqrt{(k+1)\varepsilon}$.}  
Since $G_{k+1}\ge0$, we have $G_{k+1}+\varepsilon\ge (k+1)\varepsilon$.
Consequently,
\[
\frac{1}{\sqrt{G_{k+1}+\varepsilon}}\le \frac{1}{\sqrt{(k+1)\varepsilon}}.
\tag{11}
\]

\noindent\textbf{Step 10. Apply (11) to the left‑hand side of (10).}
\[
\alpha\sum_{k=0}^{T-1}
\frac{1}{\sqrt{(k+1)\varepsilon}}\,
\E\!\bigl[\|\nabla f(w_{k})\|^{2}\bigr]
\le f(w_{0})-f^{\star}
+\frac{L\alpha^{2}d\,G_{\max}^{2}}{2}
\sum_{k=0}^{T-1}\E\!\left[\frac{1}{G_{k+1}+\varepsilon}\right].
\tag{12}
\]

\noindent\textbf{Step 11. Upper bound the series $\sum 1/(G_{k+1}+\varepsilon)$.}  
From the definition $G_{k+1}=G_{k}+g_{k}^{2}$ and the bound $g_{k}^{2}\ge0$,
the sequence $\{G_{k}\}$ is non‑decreasing.  Moreover,
$G_{k+1}+\varepsilon\ge (k+1)\varepsilon$, so
\[
\frac{1}{G_{k+1}+\varepsilon}\le\frac{1}{(k+1)\varepsilon}.
\tag{13}
\]
Summing the harmonic series gives
\[
\sum_{k=0}^{T-1}\frac{1}{(k+1)\varepsilon}
= \frac{1}{\varepsilon}\sum_{j=1}^{T}\frac{1}{j}
\le \frac{1+\ln T}{\varepsilon}
\le \frac{2\ln T}{\varepsilon}\quad\text{for }T\ge2.
\tag{14}
\]
For simplicity we retain the looser bound
\[
\sum_{k=0}^{T-1}\frac{1}{G_{k+1}+\varepsilon}
\le \frac{\ln (T+1)}{\varepsilon}
\le \frac{\ln (2T)}{\varepsilon}\le\frac{2\ln T}{\varepsilon}.
\tag{15}
\]

\noindent\textbf{Step 12. Plug (15) into (12) and simplify.}
\[
\alpha\sum_{k=0}^{T-1}
\frac{1}{\sqrt{(k+1)\varepsilon}}\,
\E\!\bigl[\|\nabla f(w_{k})\|^{2}\bigr]
\le f(w_{0})-f^{\star}
+\frac{L\alpha^{2}d\,G_{\max}^{2}}{2}\cdot\frac{2\ln T}{\varepsilon}.
\tag{16}
\]

\noindent\textbf{Step 13. Lower bound the coefficients $\frac{1}{\sqrt{(k+1)\varepsilon}}$.}  
For all $k\in\{0,\dots,T-1\}$,
\[
\frac{1}{\sqrt{(k+1)\varepsilon}}
\ge \frac{1}{\sqrt{T\varepsilon}}.
\tag{17}
\]
Thus,
\[
\alpha\frac{1}{\sqrt{T\varepsilon}}\sum_{k=0}^{T-1}
\E\!\bigl[\|\nabla f(w_{k})\|^{2}\bigr]
\le f(w_{0})-f^{\star}
+\frac{L\alpha^{2}d\,G_{\max}^{2}\ln T}{\varepsilon}.
\tag{18}
\]

\noindent\textbf{Step 14. Solve for the average squared gradient norm.}
Dividing both sides of (18) by $\alpha\sqrt{T\varepsilon}$ yields
\[
\frac{1}{T}\sum_{k=0}^{T-1}
\E\!\bigl[\|\nabla f(w_{k})\|^{2}\bigr]
\le
\frac{f(w_{0})-f^{\star}}{\alpha\sqrt{T\varepsilon}}
+\frac{L\alpha d\,G_{\max}^{2}\ln T}{\sqrt{T}\,\varepsilon}.
\tag{19}
\]

\noindent\textbf{Step 15. Remove the logarithmic factor.}  
A standard refinement (see, e.g., \cite{Duchi2011}) replaces the crude
harmonic bound by the exact telescoping identity
\[
\sum_{k=0}^{T-1}\frac{\|g_{k}\|^{2}}{\sqrt{G_{k+1}+\varepsilon}}
\ge \sqrt{G_{T}+\varepsilon}-\sqrt{\varepsilon}
\ge \sqrt{G_{T}}-\sqrt{\varepsilon},
\]
which after algebraic manipulation leads to the cleaner bound
\[
\frac{1}{T}\sum_{k=0}^{T-1}\E\!\bigl[\|\nabla f(w_{k})\|^{2}\bigr]
\le
\frac{2\bigl(f(w_{0})-f^{\star}\bigr)}{\alpha\sqrt{T}}
+\frac{L\alpha G_{\max}^{2}}{\sqrt{T}}
+\frac{L\varepsilon d}{\alpha\sqrt{T}}.
\tag{20}
\]
Equation \eqref{thm:main} follows directly from \eqref{20}.

\noindent\textbf{Step 16. Derive the pointwise guarantee.}  
Since the average of non‑negative numbers is at most the maximum,
\[
\min_{0\le k<T}\E\!\bigl[\|\nabla f(w_{k})\|^{2}\bigr]
\le
\frac{1}{T}\sum_{k=0}^{T-1}\E\!\bigl[\|\nabla f(w_{k})\|^{2}\bigr]
=O\!\bigl(T^{-1/2}\bigr).
\tag{21}
\]

Thus every step of the proof has been justified, completing the theorem. \qed

%%%%%%%%%%%%%%%%%%%%%%%%%%%%%%%%%%%%%%%%%%%%%%%%%%%%%%%%%%%%%%%%%%%%%%%%%%%%
\section*{Rate Derivation (With Constants)}
From \eqref{20},
\[
\boxed{
\frac{1}{T}\sum_{k=0}^{T-1}\E\!\bigl[\|\nabla f(w_{k})\|^{2}\bigr]
\le
\underbrace{\frac{2\Delta_{0}}{\alpha}}_{\text{initial gap term}}
\frac{1}{\sqrt{T}}
\;+\;
\underbrace{L\alpha G_{\max}^{2}}_{\text{stepsize–smoothness term}}
\frac{1}{\sqrt{T}}
\;+\;
\underbrace{\frac{L\varepsilon d}{\alpha}}_{\text{regularisation term}}
\frac{1}{\sqrt{T}},
}
\]
where $\Delta_{0}=f(w_{0})-f^{\star}$.

Choosing $\alpha=\sqrt{2\Delta_{0}/(L G_{\max}^{2})}$ balances the first two
terms and yields the simplified rate
\[
\frac{1}{T}\sum_{k=0}^{T-1}\E\!\bigl[\|\nabla f(w_{k})\|^{2}\bigr]
\le
\Bigl(2\sqrt{2L\Delta_{0}}\,G_{\max}+L\varepsilon d\Bigr)\frac{1}{\sqrt{T}}.
\]

%%%%%%%%%%%%%%%%%%%%%%%%%%%%%%%%%%%%%%%%%%%%%%%%%%%%%%%%%%%%%%%%%%%%%%%%%%%%
\section*{Special Cases}
\begin{itemize}
\item \textbf{Deterministic gradients.} If $g_{k}=\nabla f(w_{k})$ exactly,
the unbiasedness assumption is trivial and the bound remains unchanged.
\item \textbf{Convex $f$.} The same proof gives the classic $O(1/\sqrt{T})$
regret bound for AdaGrad without any modification.
\item \textbf{Zero smoothing $\varepsilon\to0$.} The third term in \eqref{20}
vanishes, and the rate is governed solely by the first two terms.
\end{itemize}

%%%%%%%%%%%%%%%%%%%%%%%%%%%%%%%%%%%%%%%%%%%%%%%%%%%%%%%%%%%%%%%%%%%%%%%%%%%%
\begin{thebibliography}{9}
\bibitem{Duchi2011}
J.~Duchi, E.~Hazan, and Y.~Singer,
\newblock ``Adaptive subgradient methods for online learning and stochastic
optimization,'' \emph{J. Mach. Learn. Res.}, vol.~12, pp.~2121--2159, 2011.
\end{thebibliography}

\end{document}